\section{Asymmetric grids and integrated spectral projections}\label{tra:asymmetric}

All the preceding grids have comparable side lengths. We now change the aspect ratio itself. Short cycles in the resulting array and fixed-size coordinate chunks supply estimates adapted to unequal child scales.

A split into a short coordinate block and a long block produces a different trace recursion. Sparse partitions are first controlled by interpolation on an enlarged tuple space, with fixed-size coordinate chunks. For the dense estimate, the associated bipartite graph has few edges belonging to a parallel pair or a four-cycle. Those edges account for the entropy error, and a support cutoff in each spectral integral pays for the number of local partition types.

Write $n=2^e$ in this section, and let $T_n$ be one sweep. Schatten norms are unnormalized, and $d_\lambda$ is the dimension of the irreducible indexed by $\lambda\vdash n$. The operator is a product of subgroup projections and has a strict norm gap off constants at each fixed size by Lemma~\ref{lem:finitegap}.
\label{tra:asymmetric:strip-interpolation}
We use the strip interpolation established in the proof of
Lemma~\ref{lem:positive-power-comparison}. For a bounded analytic
matrix family $F$ on $0\le\Re z\le1$, continuous on the boundary,
\[
 \|F(iu)\|_{\rm op}\le1,\qquad \|F(1+iu)\|_{S^q}\le B
 \quad\Longrightarrow\quad
 \|F(\theta)\|_{S^{q/\theta}}^{q/\theta}\le B^q.
\]
Below $q$ is $2$ or $4$. Complex powers are zero on their kernels,
as in that proof.
\subsection{Fixed-size coordinate chunks on enlarged tuples}
\begin{proposition}\label{tra:asymmetric:sparse}
\label{tra:asymmetric:sparse-statement}
\textbf{Long first row.} For all sufficiently large powers of two \(n\), if
\[
                     1\le k=n-\lambda_1\le n^{.995},
\]
there is a bound
\begin{equation}
               d_\lambda\,{\rm Tr}|T_n(\lambda)|^p\le n^{-10k}
                       \qquad (p\ge p_*)
               \label{tra:asymmetric:eq1}
\end{equation}
with \(p_*<\infty\) absolute.
\end{proposition}
\begin{proof}\label{tra:asymmetric:sparse-proof}
\label{tra:asymmetric:macro-coordinate-interpolation}
Use the permutation representation on ordered \(K\)-tuples of distinct positions, where \(K=\lfloor k n^{.001}\rfloor\); thus \(K/n\le n^{-.004}\). Split the coordinate sequence into chunks of lengths between \(\ell\) and \(2\ell\), where \(\ell\) will just be a sufficiently large \emph{constant}, and \(e\) can be taken sufficiently large. So the positions are tuples of macro coordinates, ranging in sets of sizes \(L_j\), each between \(2^\ell\) and \(2^{2\ell}\). A sweep updates these macro coordinates in sequence, doing sweeps \(T_{L_j}\) independently in every fibre of coordinate \(j\), by which we mean with all other macro coordinates fixed. For the ordered \(K\)-tuples this means the operator of a chunk is block diagonal, with a block specified by the allocation of tuple members to fibres. On such a block it is a tensor product: for a fibre with \(t\) members use \(T_{L_j}\) acting on ordered \(t\)-tuples (using fixed ordering of the members in this factor).

For any such local operator \(T\) write \(T=V|T|\) with unitary polar factor, and replace it by
\[
                            T(z)=V |T|^{Qz/2}.
\]
Apply these replacements in every block, and let \(F(z)\) be the product of the resulting chunk operators in the sweep order. We can take \(Q>2\) a sufficiently large constant depending on \(\ell\). On the first line of the strip the operator norms are at most 1. On the second line \(T(z)\) is arbitrarily close to the matrix averaging uniformly on its local tuple space, since \(T\) preserves constants with all other singular values strictly less than 1 by the generation observation; only finitely many local sizes are used. For \(t=0,1\) we have the averaging exactly on this line (a sweep gives a single point all fresh fair bits). In particular we take \(Q\) so that on this line, with \(L=L_j\), each entry of \(T(z)\) in absolute value is at most
\begin{equation}
                \frac{1+\delta_L(t)}{(L)_t},\qquad
       \delta_L(0)=\delta_L(1)=0,\quad \delta_L(t)=L^{-2}\ (t\ge2),
               \label{tra:asymmetric:eq2}
\end{equation}
with \((L)_t\) a falling factorial.

Between any given input and output ordered tuple, there is at most one trajectory of tuples through the chunk product: at each boundary of chunks, the already processed macro coordinates of each point must be its output ones and the others its input ones. A valid such trajectory, under \emph{ideal} routing doing full uniform permutations independently in each fibre at each chunk, has probability given by multiplying \(1/(L)_t\) for its fibres. Thus \eqref{tra:asymmetric:eq2} bounds the sum of squared entries of \(F(z)\) by
\begin{equation}
     (n)_K\ \mathbb E\prod_{\text{fibres over chunks}}
                     \frac{(1+\delta_L(t))^2}{(L)_t}
         \ \le\ \mathbb E\prod_{\text{fibres over chunks}} W_L(t),
         \qquad W_L(t)= (1+\delta_L(t))^2\frac{L^t}{(L)_t}.
                 \label{tra:asymmetric:eq3}
\end{equation}
Here we sample the starting \(K\)-tuple uniformly and do the ideal routing, \(t\) referring to the occupancy by those routed points; we used \(\prod_j L_j=n\).

We detail why clustering of the \(K\) points in this estimate is harmless. Condition on the full ideal routing of all \(n\) points, let \(X\) be the uniform size \(K\) sample of input points, and use the graph on \(n\) points connecting any two that visit the same fibre at any chunk. Maximum degree is \(O_\ell(\log n)\). Consider a component \(Y\) of the induced graph on the sample, of size \(v\). Over all chunks its fibre occupancies satisfy
\begin{equation}
                              \sum t\log t\le v\log v .
                      \label{tra:asymmetric:eq4}
\end{equation}
In fact, group the inputs by their not-yet-processed coordinates: form the tree starting with all inputs and splitting by the input coordinate of the last chunk, then the next-to-last, and so on. Take any node specified by coordinates later than \(j\), split into children by coordinate \(j\). The points of \(Y\) here (say of number \(v_0>0\)) when routed up to chunk \(j\) are distributed over fibres still within the specified later coordinates, and the \(t\) in one fibre come from \(t\) different children because their \(j\)-coordinates are distinct and not yet processed. Hence the sum of \(t\log t\) over these fibres is at most \(v_0\) times the entropy of the child proportions: average the distributions over children given by each fibre (each distribution uniform on \(t\) children with weight \(t/v_0\)) and use concavity. Summing the bounds at nodes telescopes, giving \eqref{tra:asymmetric:eq4}. There is no shared occupied fibre of two different components.

For \(t\ge 2\),
\[
                         \log W_L(t)\le \frac{C}{\log L}\,t\log t
\]
with \(C\) absolute. Up to \(\sqrt L\) use \(\log(L^t/(L)_t)=O(t^2/L)\); above that, it suffices that this log is at most \(t\), using the average over the full increasing sequence \(\log(L/(L-i)),\,0\le i<L\). Choose \(\ell\) so large that by this and \eqref{tra:asymmetric:eq4} the product in \eqref{tra:asymmetric:eq3} is bounded by the product over nonsingleton components of
\(v^{\zeta v}\), \(\zeta=.0005\).

With \(\Delta=O_\ell(\log n)\) a degree bound (at least 1), and \(u=K/n\), the expectation of this last product is at most
\[
                  \exp\left(\sum_{v=2}^K n\Delta^{2(v-1)} u^v v^{\zeta v}\right).
\]
Indeed we can sum products over collections of disjoint connected sets included in \(X\), charging the components in particular; the probability cost for such a collection is at most \(u\) to the power of total size. The count of connected sets uses a spanning tree walk of length \(2(v-1)\); allowing arbitrary collections with the product cost gives the exponential bound. The sum in this exponential is \(O(n\Delta^2 u^2)\): for \(v\ge3\) compare after factoring \(n\Delta^2 u^2\) to powers of \(\Delta^2 u n^{3\zeta}\), which goes to zero. Thus \eqref{tra:asymmetric:eq3} is at most \(\exp(O(k))\), since \(\Delta^2 K^2/n=O(k)\) for our \(k,K\).

The enlarged tuple space has now supplied a Schatten budget
$\exp(O(k))$. Its large branching multiplicity will convert that budget
into decay on each first-row level. We now interpolate. We have \(\|F(1+iy)\|_2^2\le\exp(O(k))\), whereas \(F(2/Q)\) is the sweep operator on the \(K\)-tuple space, so its Schatten \(Q\)-norm to power \(Q\) has the same upper bound. By branching and Frobenius reciprocity for the pointwise stabilizer of \(K\) positions, the multiplicity of \(\lambda\) in this space is the number of standard tableaux of the skew shape obtained by removing the row of length \(n-K\) from \(\lambda\). Since \(\lambda_2\le n-K\) here (for large enough \(n\)), there are at least \(\binom Kk\) such tableaux: the tail of size \(k\) is separate from the remaining \(K-k\) entries of the first row. The Schatten estimate with this multiplicity thus gives, say,
\[
                         \|T_n(\lambda)\|_{\rm op}^Q\le n^{-.0004 k}
\]
for all sufficiently large \(n\) (use \(K/k\ge n^{.001}/2\)). As \(d_\lambda\le n^k\) (it occurs even on ordered \(k\)-tuples), increasing to a suitably large constant power \(p_*\) proves \eqref{tra:asymmetric:eq1}.
\end{proof}
\begin{lemma}\label{tra:asymmetric:dimension}
\label{tra:asymmetric:dimension-and-annihilation}
We record two further partition details. If \(\lambda\) has length greater than \(n/2\), we have \(T_n(\lambda)=0\). Indeed there are no invariants under even one full pair layer subgroup. This also follows immediately in the Specht module realization: the irrep is spanned by column-antisymmetrized vectors (polytabloids) from tableaux, and the column of length greater than \(n/2\) must contain both endpoints of some pair, so averaging kills such vectors. For all sufficiently large \(n\), for the partitions of length at most \(n/2\) with \(k>n^{.995}\), we have
\begin{equation}
                            \log d_\lambda\ \ge\ n^{.995}.
                      \label{tra:asymmetric:eq5}
\end{equation}
In fact if the first row has length \(r\ge n/10\), keep a subshape of the diagram with this row and \(t=\lfloor n^{.995}/2\rfloor\) boxes below. Even on this shape we get at least \(\binom rt\) tableaux: after the first \(t\) boxes of the first row are filled we can interleave the rest of that row and a standard filling of the lower boxes. They extend to the whole shape. This count suffices for \eqref{tra:asymmetric:eq5}. We can do the same with the first column if it has length at least \(n/10\) (there are enough boxes outside the column by the length condition). If neither holds, \eqref{tra:asymmetric:eq5} follows by the hook length formula with all hooks at most \(n/5\).
\end{lemma}
\subsection{Short-cycle entropy on an asymmetric grid}
\label{tra:asymmetric:grid-sizes}
For the regular-representation part of the proof, take, with \(n=2^e\) large,
\[
             a=2^{\lfloor e/10\rfloor},\qquad b=n/a.
\]
\begin{proposition}\label{tra:asymmetric:compatibility}
\label{tra:asymmetric:compatibility-statement}
Arrange the points in \(b\) rows of \(a\) cells each. We need a probability bound, which we prove first. Independently in each row \(i\), assign the labels \(h_{ij}\in[a]\) by a permutation, with an arbitrary law of density at most \(D_i\) relative to uniform. Independently each column \(j\) has labels \(l_{ij}\in[b]\) likewise assigned by a permutation, say with density at most \(D'_j\), also independent across columns and from the row assignments. Put \(D_H=\prod_i D_i,\ D_L=\prod_j D'_j\). Then, with an absolute constant \(C\), the probability that the pairs \((h_{ij},l_{ij})\) are all distinct is at most
\begin{equation}
                \exp\left(-n+C n^{.98}
                        +\frac{C}{\log n}\log(D_H D_L)\right).
                     \label{tra:asymmetric:eq6}
\end{equation}
\end{proposition}
\begin{proof}\label{tra:asymmetric:compatibility-proof}
\label{tra:asymmetric:short-cycle-and-matching-entropy}
For the proof, the \(l\)-array gives a bipartite multigraph on sides of size \(b\), with one edge from \(i\) to \(l_{ij}\) per cell; it has degree \(a\). Call an edge bad if contained in a parallel pair or in a 4-cycle. Write \(L_{\rm bad}\) for the number of bad edges (cells).

Here is a useful moment bound for these edges under the uniform independent column assignments \(\mathcal U\):
\begin{equation}
        \log\mathbb E_{\mathcal U}\exp(c(\log n)L_{\rm bad})
                                      = O(n^{.9})
                       \label{tra:asymmetric:eq7}
\end{equation}
for some absolute \(c>0\). Compare any law \(\xi\) of column assignments to \(\mathcal U\) by KL divergence (relative entropy). Use the chain rule revealing one cell at a time in an order given by independent uniform priorities. For a current cell, if at least \(b/3\) values are still available in its column, test the event that its new edge closes a parallel pair or a 4-cycle with edges already visible; otherwise use an empty event as test. Under the \(\mathcal U\)-conditional distribution, chance of the test event is \(O(a^3/b)=O(n^{-.6})\); the possible targets from paths of length 1 or 3 are bounded in number by \(a+a^3\). Thus each conditional KL divergence is bounded below by
\[
           .25(\log n)\Pr_{\xi}(\text{test event}\mid\text{past})-O(n^{-.35}),
\]
where the order is given (use the entropy inequality with \(.25\log n\) times the indicator as test function). For a given completed array, each bad edge will score the event with at least absolute positive probability over order. Indeed fix up to three witnessing other edges; with probability bounded below the cell priority is in \((.3,.4)\) and these witnesses precede it, and at least \(b/3\) column choices are still available. For the last condition one can use just concentration of the number of other cells of that column seen, with at most three additional cells forced earlier. Averaging and summing, the divergence of \(\xi\) from \(\mathcal U\) is at least \(c(\log n)\mathbb E_\xi L_{\rm bad}-O(n^{.9})\), giving \eqref{tra:asymmetric:eq7}, for example by taking \(\xi\) the exponential tilt. Hence for our possibly nonuniform column laws, by Hölder (with exponent \(c\log n\) on \(\exp L_{\rm bad}\), and using density at most \(D_L\) of the product law),
\begin{equation}
             \log\mathbb E\exp L_{\rm bad}
                         \le O(n^{.9}) + \frac{\log D_L}{c\log n}.
                       \label{tra:asymmetric:eq8}
\end{equation}

The short-cycle estimate controls the column geometry. We next
condition on that geometry and expose the row permutations; the
restriction to edges belonging to neither a parallel pair nor a
four-cycle makes the blocking rows distinct.

Fix columns and upper bound the conditional probability \(p\) of distinct pairs, assuming it is positive. The conditional law of row assignments \textbf{given} distinct pairs has divergence \(-\log p\) from the product of the original row laws. Reveal rows in a random order, again using independent uniform priorities. At row \(i\), a conditional divergence term is at least minus log of the original row-law mass of candidates fitting with previously revealed rows (no reused pairs), since the conditional law is supported on such candidates.

To bound these masses, fix any full realization from the law given distinct pairs. In row \(i\) let \(g_i\) be the number of cells whose column-array edge is not bad. For each such cell \(ij\) where a candidate for the row changes the \(h\)-label, the new pair (new \(h\), \(l_{ij}\)) occurs in the full realization and must occur on some other row (as \(l_{ij}\) doesn't repeat within row \(i\)). Moreover all the blocking rows identified this way for the candidate are distinct: otherwise two such \(l\) values connect row \(i\) and another row by a 4-cycle. If \(v_i\) is the priority of \(i\), put
\[
                  x=\min\{-\log(1-v_i), .25\log a\}.
\]
Thus for a given candidate matching the realized row permutation at \(M\) cells (total matches), the probability over remaining priorities of fitting is at most
\(\exp(-x g_i+x M)\), since none of the blocking rows can have come before it. So the expected log of the mass fitting, given the realization and \(v_i\), by Jensen is at most
\[
                     -x g_i + \log\mathbb E_{\nu_i}\exp(xM)
\]
where \(\nu_i\) is the original candidate law on row \(i\). For uniform permutations the log of the moment \(\mathbb E\exp(tM)\), \(t\ge 0\), is at most \(\exp(t)-1\), by counting subsets of matches. Thus by Hölder using exponent \(.5(\log a)/x\) on the moment function for \(x>0\), we get
\[
                    \log\mathbb E_{\nu_i}\exp(xM)
                      \le \frac{2x}{\log a}(\log D_i + a^{.5}).
\]
We integrate these log mass bounds keeping in mind \(\mathbb E x=1-a^{-1/4}\). They apply uniformly when histories are taken from the full law conditioned on distinct pairs (which can be sampled independently of the priorities). So by the divergence bound for \(-\log p\),
\[
           \log p
                 \le -(1-a^{-1/4})(n-L_{\rm bad})
                       +\frac{2}{\log a}\left(\log D_H+b a^{.5}\right)
                 \le -n+ L_{\rm bad} + O(n^{.98})
                                +\frac{C}{\log n}\log D_H.
\]
This along with \eqref{tra:asymmetric:eq8} proves \eqref{tra:asymmetric:eq6}.
\end{proof}
\begin{corollary}\label{tra:asymmetric:projection-crossing}
\label{tra:asymmetric:projection-crossing-statement}
We apply the estimate to products of group projections. Use \(R\) and \(Y\) for the subgroups of \(S_n\) permuting within rows and within columns, respectively. For each row of size \(a\), let \(P_i\) be an orthogonal projection in its group algebra supported just on one irrep \(\mu_i\), of rank \(r_i>0\) there (so acting on the regular representation, this block action is repeated with multiplicity \(d_{\mu_i}\)). Put \(D_i=d_{\mu_i}r_i\). Likewise take projections for the columns in their group algebras with \(D'_j\) defined similarly. Write \(P_R,P_Y\) for the two product projections, viewed also in the group algebra of \(S_n\). On its regular representation,
\begin{equation}
           \|P_Y P_R\|_4^4
                  \le \exp(C n^{.98})\,
                         \left(\prod_i D_i \prod_j D'_j\right)^{1+C/\log n}
                        \label{tra:asymmetric:eq9}
\end{equation}
for an absolute \(C\).
\end{corollary}
\begin{proof}\label{tra:asymmetric:projection-crossing-proof}
\label{tra:asymmetric:projection-crossing-proof-details}
Apply Lemma~\ref{lem:coefficient-compatibility} with the line
projections as its positive operators. Their regular ranks and squared
masses are $D_i=d_{\mu_i}r_i$ and $D'_j$, so the normalized squared
coefficient densities have these same caps. At an intermediate cell
between row action $r$ and column action $y$, the column of origin is
given by $r^{-1}$ and the row of destination by $y$. Thus opposite-order
routing is exactly the distinct-pair event in
Proposition~\ref{tra:asymmetric:compatibility}; inversion preserves
the coefficient caps. The lemma and \eqref{tra:asymmetric:eq6} give
\[
 \|P_YP_R\|_4^4
 \le\frac{n!}{|R||Y|}
       \left(\prod_iD_i\prod_jD'_j\right)
       \exp\left\{-n+Cn^{.98}
          +\frac{C}{\log n}\log\left(\prod_iD_i\prod_jD'_j\right)\right\}.
\]
Finally $n!/(|R||Y|)=n!/((a!)^b(b!)^a)\le e^{n+O(\log n)}$,
which proves \eqref{tra:asymmetric:eq9}.
\end{proof}
\subsection{Spectral integrals with a dimension cutoff}
\begin{theorem}\label{tra:asymmetric:trace}
\label{tra:asymmetric:trace-statement}
We prove that for each sweep size there is \(p_n\ge 4\), all bounded above by an absolute constant, such that
\begin{equation}
                     d_\lambda\,{\rm Tr}|T_n(\lambda)|^{p_n}
                            \ \le\ \exp(n^{.99})
                     \qquad\text{for all }\lambda.
                          \label{tra:asymmetric:eq10}
\end{equation}
For all sufficiently large $n$, the proof also bounds the complete regular Schatten power by $\exp(n^{.99})$.
\end{theorem}
\begin{proof}\label{tra:asymmetric:trace-proof}
\label{tra:asymmetric:spectral-integral-recursion}
Let \(B\ge 1\) be a fixed constant sufficiently large for \eqref{tra:asymmetric:eq9}, so we can use exponent \(\Gamma_n=1+B/\log n\) in that bound. We use a sufficiently large absolute cutoff \(n_0\), which can be taken above all fixed thresholds as needed in the proof below (in particular the estimates requiring sufficiently large sizes in the induction step will not depend on the base choice of power). Choose \(p_{\rm base}\ge \max\{4,p_*\}\), also sufficiently large to make the left side of \eqref{tra:asymmetric:eq10} with \(p_{\rm base}\) at most 1 for every size up to the cutoff; this is possible by strict contraction on nontrivial sectors. Use \(p_{\rm base}\) up to the cutoff. Above, with \(a,b\) from the grid estimate, set
\[
                  p_n=\left(1+\frac{B+1}{\log n}\right)\max(p_a,p_b).
\]
This is bounded absolutely: down any recursion path above the cutoff, the logs of sizes decrease by at least a constant factor (i.e. next at most, say, .95 times the current log), so the sum of their reciprocals there is bounded.

Split the coordinates into the first \(\log_2 a\) and the remaining \(\log_2 b\), arranging the wires so that \(T_n=T_Y T_R\) with \(T_R\) a product over rows of copies of \(T_a\), and \(T_Y\) likewise over columns with \(T_b\). Use the ambient regular representation. Put \(X=|T_Y|,\ Z=|T_R^*|\). We can drop outer unitary polar factors, and use
\begin{equation}
               \|T_n\|_{p_n}^{p_n}
                     \le \big\|X^{p_n/4} Z^{p_n/4}\big\|_4^4 .
                           \label{tra:asymmetric:eq11}
\end{equation}
This is Lemma~\ref{lem:positive-power-comparison} with $P=p_n$ and $q=4$.

The positive powers on the right factor as corresponding products commuting within the line subgroups (row factors amongst themselves, and likewise columns), since the operations factor in the respective subgroup algebras and lift to the ambient group algebra. Write \(p=p_n\). On a line of size \(m\) (\(a\) or \(b\)), according to orientation, use the spectral projections within irrep \(\mu\) of that group's \(|T_m|\) or \(|T_m^*|\); write \(E_{\mu,t}\) for the projection for singular values at least \(t>0\) (zero on other irreps), of rank \(r_\mu(t)\) on the irrep. The line power for \eqref{tra:asymmetric:eq11} is
\[
            \sum_\mu \int_0^1 \frac{p}{4}\,t^{p/4-1} E_{\mu,t}\,dt
\]
lifted as an element of the corresponding group algebra. Expand in all lines and use the triangle inequality and \eqref{tra:asymmetric:eq9} on the right of \eqref{tra:asymmetric:eq11}. This gives
\begin{equation}
 \|T_n\|_{p}^{p}
   \le \exp(C n^{.98})
        \prod_{\text{lines}}
        \left(\sum_\mu \int_0^1
             \frac{p}{4}\,t^{p/4-1}
             \big(d_\mu r_\mu(t)\big)^{\Gamma_n/4}\,dt\right)^4 .
                  \label{tra:asymmetric:eq12}
\end{equation}

The spectral support cutoff supplies a negative power of a large child
dimension. This compensates for the number of local partition types:
a slight surplus in the parent exponent becomes a dimension saving
because the integral ends before $t=1$. We show that each integral
in \eqref{tra:asymmetric:eq12} is $O(\log n)$. Write \(W=d_\mu {\rm Tr}|T_m(\mu)|^{p_m}\). Then
\[
         d_\mu r_\mu(t)\le W t^{-p_m},\qquad
         r_\mu(t)=0\ \text{ if }\ t>(W/d_\mu)^{1/p_m}.
\]
If \(W\le1\) then simply integrating with the first bound gives at most \(p/(p-\Gamma_n p_m)=O(\log n)\). This applies for sizes up to the cutoff, and also in larger sizes for the trivial irrep, for the cases from \eqref{tra:asymmetric:eq1}, and for the sectors killed by the pair layer. For the remaining cases, we can use \(W\le H=\exp(m^{.99})\) inductively and have \(\log d_\mu\ge m^{.995}\) by \eqref{tra:asymmetric:eq5}. Using the support bound as well gives at most
\[
          \frac{p}{p-\Gamma_n p_m}
                   H^{\Gamma_n/4} (H/d_\mu)^{(p/p_m-\Gamma_n)/4}.
\]
Since \(p/p_m-\Gamma_n\ge 1/\log n\), this is \(O(\log n)\) as claimed (for \(m\) above the cutoff), using \(m^{.995}\gg m^{.99}\log n\) for large sizes here, as \(\log n=O(\log m)\).

The count of partitions of \(m\) is at most \(\exp(O(\sqrt m\log(m+1)))\), for example by describing rows and columns out from the diagonal. Therefore \eqref{tra:asymmetric:eq12} with \(b\) lines of size \(a\) and \(a\) lines of size \(b\) yields
\[
                              \|T_n\|_p^p\le \exp(n^{.99}),
\]
since \(b\sqrt{a}=O(n^{.95})\), \(a\sqrt{b}=O(n^{.55})\). All the estimates needed here for sufficiently large \(n\) have constants independent of \(p_{\rm base}\), so a fixed cutoff as specified suffices. We have bounded the regular-representation sum, proving in particular \eqref{tra:asymmetric:eq10} and closing the induction.
\end{proof}
\begin{corollary}[Vanishing regular remainder]\label{tra:asymmetric:mixing}
Let $p^*=\sup_n p_n<\infty$ be the bound obtained above and let
$\Pi$ be uniform averaging. For every integer $j\ge2p^*$,
$\|T_n^j-\Pi\|_{\mathrm{HS}}^2\longrightarrow0$ as $n\to\infty$.
\end{corollary}
\begin{proof}\label{tra:asymmetric:mixing-proof}
\label{tra:asymmetric:fourier-completion}
Finally take a constant integer number of sweeps \(j\) at least twice the upper bound on the \(p_n\)'s. For nontrivial sectors with \(1\le k\le n^{.995}\), \eqref{tra:asymmetric:eq1} shows for large sizes
\[
            d_\lambda \| T_n(\lambda)^j \|_2^2
                   \le d_\lambda^2 \| T_n(\lambda)\|_{\rm op}^{2j}
                   \le d_\lambda^2 (n^{-10k}/d_\lambda)^4
\]
where the first norm after the dimension factor is Hilbert-Schmidt. This is negligible even summed (there are at most \(2^k\) partitions for a given \(k\)). For remaining sectors not already killed we similarly have by \eqref{tra:asymmetric:eq10}
\[
          d_\lambda \| T_n(\lambda)^j\|_2^2
                  \le d_\lambda^2 (\exp(n^{.99})/d_\lambda)^4
\]
which by \eqref{tra:asymmetric:eq5} is again negligible summed over partitions. Thus the squared Hilbert-Schmidt norm in the regular representation of \(T_n^j\) minus uniform averaging tends to zero. The identification with the squared $L^2$ distance of the permutation
density is the one in Proposition~\ref{e:nonnormal-moment-conversion}.
All powers above are controlled by operator submultiplicativity; no
normality of $T_n$ is required.
\end{proof}
